\documentclass[11pt,aspectratio=169]{ctexbeamer}

\usetheme{default}
\setbeamertemplate{navigation symbols}{}

\usepackage{amsmath}
\usepackage{booktabs}

% Projector-safe palette: high contrast on common lecture-room projectors.
\definecolor{DeepBlue}{RGB}{11,58,102}
\definecolor{SoftInk}{RGB}{31,41,51}
\definecolor{PaperBg}{RGB}{255,255,255}
\definecolor{PanelBlue}{RGB}{226,237,247}

\setbeamercolor{normal text}{fg=SoftInk,bg=PaperBg}
\setbeamercolor{structure}{fg=DeepBlue}
\setbeamercolor{frametitle}{fg=DeepBlue,bg=PaperBg}
\setbeamercolor{title}{fg=DeepBlue}
\setbeamercolor{block title}{fg=DeepBlue,bg=PanelBlue}
\setbeamercolor{block body}{fg=SoftInk,bg=white}

\newcommand{\Core}[1]{
  \begin{block}{核心结论}
    #1
  \end{block}
}

\newcommand{\LectureMeta}[2]{
  \title{大模型推理系统与实践}
  \subtitle{#1}
  \author{张书豪}
  \institute{华中科技大学计算机学院}
  \date{2027 · #2}
}

\newcommand{\PaperRef}[3]{
  \item \textbf{#1}（#2）\hfill\href{#3}{原文链接}
}

\newcommand{\ResearchSources}[1]{
  \begin{frame}{本讲参考资料}
    \normalsize
    \begin{itemize}
      #1
    \end{itemize}
    \vspace{0.3em}
    \textcolor{SoftInk}{课件中的论文数字均指作者在特定模型、硬件和工作负载下报告的结果，不代表无条件性能。}
  \end{frame}
}

\newcommand{\InternalCaseNote}{
  \small\textcolor{SoftInk}{说明：本页为课程研究案例或教学假设，不代表已发表论文结论；只有在实验包给出模型、硬件、工作负载和原始记录时，数字才可作为结果引用。}
}

% 14 讲细分版的统一参考资料入口。论文数字的解释边界见课程内容基线。

\newcommand{\LectureOneSources}{\ResearchSources{
  \PaperRef{vLLM: Easy, Fast, and Cheap LLM Serving with PagedAttention}{SOSP 2023}{https://doi.org/10.1145/3600006.3613165}
  \PaperRef{SGLang: Efficient Execution of Structured Language Model Programs}{arXiv:2312.07104}{https://arxiv.org/abs/2312.07104}
  \PaperRef{nano-vllm-hust 教学代码}{课程仓库}{https://github.com/intellistream/intro-to-llm-inference-engines}
}}

\newcommand{\LectureTwoSources}{\ResearchSources{
  \PaperRef{vLLM / PagedAttention}{SOSP 2023}{https://doi.org/10.1145/3600006.3613165}
  \PaperRef{Sarathi-Serve}{OSDI 2024}{https://www.usenix.org/conference/osdi24/presentation/agrawal}
  \PaperRef{DistServe}{OSDI 2024}{https://www.usenix.org/conference/osdi24/presentation/zhong-yinmin}
}}

\newcommand{\LectureThreeSources}{\ResearchSources{
  \PaperRef{vLLM / PagedAttention}{SOSP 2023}{https://doi.org/10.1145/3600006.3613165}
  \PaperRef{DistServe}{OSDI 2024}{https://www.usenix.org/conference/osdi24/presentation/zhong-yinmin}
  \PaperRef{Llumnix}{OSDI 2024}{https://www.usenix.org/conference/osdi24/presentation/sun-biao}
}}

\newcommand{\LectureFourSources}{\ResearchSources{
  \PaperRef{Orca}{OSDI 2022}{https://www.usenix.org/conference/osdi22/presentation/yu}
  \PaperRef{Sarathi-Serve}{OSDI 2024}{https://www.usenix.org/conference/osdi24/presentation/agrawal}
  \PaperRef{FastServe}{NSDI 2026}{https://www.usenix.org/conference/nsdi26/presentation/wu-bingyang}
}}

\newcommand{\LectureFiveSources}{\ResearchSources{
  \PaperRef{vLLM / PagedAttention}{SOSP 2023}{https://doi.org/10.1145/3600006.3613165}
  \PaperRef{vAttention}{arXiv:2405.04437}{https://arxiv.org/abs/2405.04437}
  \PaperRef{Rethinking KV Cache Compression}{MLSys 2025}{https://proceedings.mlsys.org/paper_files/paper/2025/hash/26289c647c6828e862e271ca3c490486-Abstract-Conference.html}
  \PaperRef{Mooncake}{FAST 2025}{https://www.usenix.org/conference/fast25/presentation/qin}
}}

\newcommand{\LectureSixSources}{\ResearchSources{
  \PaperRef{DistServe}{OSDI 2024}{https://www.usenix.org/conference/osdi24/presentation/zhong-yinmin}
  \PaperRef{RAGCache}{arXiv:2404.12457}{https://arxiv.org/abs/2404.12457}
  \PaperRef{Mooncake}{FAST 2025}{https://www.usenix.org/conference/fast25/presentation/qin}
  \PaperRef{NEO}{MLSys 2025}{https://proceedings.mlsys.org/paper_files/paper/2025/hash/66a026c0d17040889b50f0dfa650e5e0-Abstract-Conference.html}
}}

\newcommand{\LectureSevenSources}{\ResearchSources{
  \PaperRef{vLLM / PagedAttention}{SOSP 2023}{https://doi.org/10.1145/3600006.3613165}
  \PaperRef{SGLang}{arXiv:2312.07104}{https://arxiv.org/abs/2312.07104}
  \PaperRef{DistServe}{OSDI 2024}{https://www.usenix.org/conference/osdi24/presentation/zhong-yinmin}
  \PaperRef{Mooncake}{FAST 2025}{https://www.usenix.org/conference/fast25/presentation/qin}
}}

\newcommand{\LectureEightSources}{\ResearchSources{
  \PaperRef{FlashAttention}{NeurIPS 2022}{https://arxiv.org/abs/2205.14135}
  \PaperRef{FlexGen}{ICML 2023}{https://proceedings.mlr.press/v202/sheng23a.html}
  \PaperRef{PowerInfer}{SOSP 2024 / arXiv:2312.12456}{https://arxiv.org/abs/2312.12456}
  \PaperRef{FlashInfer}{MLSys 2025}{https://proceedings.mlsys.org/paper_files/paper/2025/hash/dbf02b21d77409a2db30e56866a8ab3a-Abstract-Conference.html}
}}

\newcommand{\LectureNineSources}{\ResearchSources{
  \PaperRef{DEEPSERVE}{USENIX ATC 2025}{https://www.usenix.org/conference/atc25/presentation/hu-junhao}
  \PaperRef{QServe}{MLSys 2025}{https://proceedings.mlsys.org/paper_files/paper/2025/hash/fbe2b2f74a2ece8070d8fb073717bda6-Abstract-Conference.html}
  \PaperRef{vLLM-HUST}{开源维护仓库}{https://github.com/vLLM-HUST/vllm-hust}
}}

\newcommand{\LectureTenSources}{\ResearchSources{
  \PaperRef{vLLM / PagedAttention}{SOSP 2023}{https://doi.org/10.1145/3600006.3613165}
  \PaperRef{Revisiting Disaggregated LLM Serving}{2025 preprint, arXiv:2601.08833}{https://arxiv.org/abs/2601.08833}
  \PaperRef{Rethinking KV Cache Compression}{MLSys 2025}{https://proceedings.mlsys.org/paper_files/paper/2025/hash/26289c647c6828e862e271ca3c490486-Abstract-Conference.html}
}}

\newcommand{\LectureElevenSources}{\ResearchSources{
  \PaperRef{MLPerf Inference}{MLCommons benchmark}{https://mlcommons.org/benchmarks/inference-datacenter/}
  \PaperRef{ACM Artifact Review and Badging}{ACM policy}{https://www.acm.org/publications/policies/artifact-review-and-badging-current}
  \PaperRef{Revisiting Disaggregated LLM Serving}{2025 preprint, arXiv:2601.08833}{https://arxiv.org/abs/2601.08833}
}}

\newcommand{\LectureTwelveSources}{\ResearchSources{
  \PaperRef{nano-vllm-hust}{课程教学代码}{https://github.com/intellistream/intro-to-llm-inference-engines}
  \PaperRef{vLLM-HUST}{HUST 维护的 vLLM fork}{https://github.com/vLLM-HUST/vllm-hust}
  \PaperRef{上游 vLLM}{开源项目}{https://github.com/vllm-project/vllm}
}}

\newcommand{\LectureThirteenSources}{\ResearchSources{
  \PaperRef{课程内容基线与实验材料}{课程仓库}{https://github.com/intellistream/intro-to-llm-inference-engines}
  \PaperRef{DEEPSERVE}{USENIX ATC 2025}{https://www.usenix.org/conference/atc25/presentation/hu-junhao}
  \PaperRef{GhostServe}{MLSys 2026}{https://proceedings.mlsys.org/paper_files/paper/2026/hash/d99e8e80a6c41e148db686918dd7eab3-Abstract-Conference.html}
}}

\newcommand{\LectureFourteenSources}{\ResearchSources{
  \PaperRef{课程完整材料与后续更新}{课程仓库}{https://github.com/intellistream/intro-to-llm-inference-engines}
  \PaperRef{SpecGen}{MLSys 2026}{https://proceedings.mlsys.org/paper_files/paper/2026/hash/66a026c0d17040889b50f0dfa650e5e0-Abstract-Conference.html}
  \PaperRef{GhostServe}{MLSys 2026}{https://proceedings.mlsys.org/paper_files/paper/2026/hash/d99e8e80a6c41e148db686918dd7eab3-Abstract-Conference.html}
}}



\LectureMeta{第 3 讲 一次推理请求的完整链路}{研究生课程课件}

\begin{document}

\begin{frame}[plain]
  \titlepage
\end{frame}

\begin{frame}{教学目标}
  \Core{本讲分解一次生成请求的完整系统链路，建立面向时序与阶段的分析视角。}
  \begin{itemize}
    \item 理解 API 到流式返回的关键阶段。
    \item 区分 prefill 与 decode 的资源特性。
    \item 找到排队点、状态生成点与资源回收点。
  \end{itemize}
\end{frame}

\begin{frame}{请求生命周期总览}
  \Core{一次请求通常要经历接入、规范化、排队、执行、返回和回收六个阶段。}
  \begin{enumerate}
    \item API 接入与参数校验
    \item tokenizer 与输入规范化
    \item 请求排队与调度
    \item prefill 生成首轮 KV
    \item decode 迭代输出 token
    \item 流式返回与资源回收
  \end{enumerate}
\end{frame}

\begin{frame}{为什么生命周期视角是基础}
  \Core{只有先把请求穿过系统的完整路径画清楚，后续任何优化才知道自己在改哪一段。}
  \begin{itemize}
    \item 生命周期视角帮助区分排队问题和执行问题。
    \item 生命周期视角帮助区分计算状态和长期状态对象。
    \item 生命周期视角帮助解释 TTFT 和 TPOT 为什么受不同阶段支配。
  \end{itemize}
\end{frame}

\begin{frame}{prefill 与 decode}
  \Core{prefill 和 decode 在资源需求、控制目标和瓶颈形式上显著不同。}
  \begin{itemize}
    \item prefill 更计算密集，直接决定首 token 延迟。
    \item decode 单轮较轻，但迭代频繁，更依赖状态驻留。
    \item 把二者混为同一任务，容易造成错误调度与资源竞争。
  \end{itemize}
\end{frame}

\begin{frame}{prefill 为什么敏感}
  \Core{prefill 直接决定首 token 何时出现，因此它既是计算问题，也是用户体验问题。}
  \begin{itemize}
    \item 长 prompt 会显著放大 prefill 开销。
    \item RAG 或工具注入会进一步增加 prompt 长度和波动。
    \item prefill 常常也是资源切分和阶段解耦的起点。
  \end{itemize}
\end{frame}

\begin{frame}{decode 为什么难管}
  \Core{decode 每轮算得不重，但它持续依赖状态驻留和低抖动调度，因此更容易受系统噪声影响。}
  \begin{itemize}
    \item 每轮都依赖已有 KV 和 batch 组织。
    \item 并发请求越多，轮次间调度和状态竞争越复杂。
    \item decode 是服务稳态控制的核心阶段。
  \end{itemize}
\end{frame}

\begin{frame}{系统中的潜在瓶颈点}
  \Core{生命周期视角的价值，在于把不同阶段的瓶颈显式区分开。}
  \begin{itemize}
    \item 队列等待与队头阻塞。
    \item prefill 对长 prompt 的容量和算力压力。
    \item decode 的状态驻留和频繁调度开销。
    \item 输出侧流式接口和回收延迟。
  \end{itemize}
\end{frame}

\begin{frame}{状态对象沿着哪条链路传播}
  \Core{一次请求在生命周期中会持续生成和消耗多类状态对象。}
  \begin{itemize}
    \item tokenizer 生成规范化输入状态。
    \item prefill 生成首轮 KV。
    \item decode 追加 KV 并消费队列与调度元数据。
    \item 返回阶段触发日志、统计和资源回收。
  \end{itemize}
\end{frame}

\begin{frame}{源码穿插：在 nano-vllm-hust 中跟踪请求路径}
  \Core{\texttt{\detokenize{nano-vllm-hust}} 保留了从请求进入到全部完成的完整主路径，足够支撑本讲的生命周期分析。}
  \begin{itemize}
    \item \texttt{\detokenize{nanovllm/engine/llm_engine.py}} 中的 \texttt{\detokenize{add_request()}} 对应请求进入系统。
    \item 同一文件中的 \texttt{\detokenize{step()}} 串联 \texttt{\detokenize{schedule()}}、模型执行和 \texttt{\detokenize{postprocess()}}。
    \item \texttt{\detokenize{generate()}} 展示了从批量请求注入到全部完成的完整外层驱动循环。
  \end{itemize}
\end{frame}

\begin{frame}{源码穿插：课堂观察重点}
  \Core{沿着 \texttt{\detokenize{LLMEngine}} 的主路径阅读源码，可以将本讲的生命周期划分与真实实现逐段对应起来。}
  \begin{itemize}
    \item 请求进入点：\texttt{\detokenize{LLMEngine.add_request()}} 中的 tokenizer 编码与 \texttt{\detokenize{Sequence}} 构造。
    \item 阶段切换点：\texttt{\detokenize{step()}} 中 \texttt{\detokenize{schedule()}} 返回的 \texttt{\detokenize{is_prefill}} 标志。
    \item 完成与回收点：\texttt{\detokenize{postprocess()}} 更新状态并在完成时输出结果。
  \end{itemize}
\end{frame}

\begin{frame}{为什么生命周期视角重要}
  \Core{很多“优化没生效”的根本原因，是优化点并不位于主导端到端时延的阶段。}
  \begin{itemize}
    \item 局部 kernel 更快，可能仍然输给排队放大。
    \item 共享执行更强，可能仍然输给 KV 驻留边界。
    \item 需要先定位主导阶段，再定义优化对象。
  \end{itemize}
\end{frame}

\begin{frame}{生命周期分析中的典型误判}
  \Core{将 GPU 利用率偏低直接视为主导问题，往往会忽略等待放大与状态竞争的真实影响。}
  \begin{itemize}
    \item 设备不满载不一定说明算力不足。
    \item 可能是请求排队、KV 碎片或重组时机不合适。
    \item 生命周期分析的作用，就是避免过早跳到算子层结论。
  \end{itemize}
\end{frame}

\begin{frame}{课堂练习}
  \Core{课堂练习可围绕请求时序图展开，以强化生命周期分析能力。}
  \begin{itemize}
    \item 标出至少 4 个潜在瓶颈点。
    \item 标出至少 3 类状态对象的生成与释放时刻。
    \item 讨论哪个阶段最可能支配 TTFT，哪个阶段最可能支配 TPOT。
  \end{itemize}
\end{frame}

\begin{frame}{讲解：为什么生命周期是后续所有机制课的前提}
  \Core{如果一次请求的阶段边界没画清楚，调度、KV、执行优化都会变成悬空讨论。}
  \begin{itemize}
    \item 生命周期把一个请求拆成进入、规范化、排队、prefill、decode、返回、回收等阶段。
    \item 不同阶段的资源特征不同，因此瓶颈也不同；TTFT 和 TPOT 往往不是由同一阶段支配。
    \item 只有先把阶段画出来，后面才知道一个优化到底改了哪里、为什么可能生效。
  \end{itemize}
\end{frame}

\begin{frame}{讲解：prefill 与 decode 为什么必须分开教}
  \Core{很多初学者把“生成”当成单一过程，但系统里最重要的区分之一就是 prefill 与 decode。}
  \begin{itemize}
    \item prefill 更像一次性大计算，直接决定首 token 何时出现。
    \item decode 更像持续迭代的稳态过程，更容易受到调度、状态驻留和抖动影响。
    \item 课程后续很多机制都是围绕这一区分展开的，例如阶段解耦、continuous batching 和状态搬运。
  \end{itemize}
\end{frame}

\begin{frame}{例子：一次多轮对话的真实链路并不止于模型前向（1）}
  \Core{请求生命周期要覆盖从输入组织到状态回收的完整过程，而不是只盯住模型执行。}
\begin{block}{进入前的准备}
    chat template、tokenizer 和 prompt 拼接本身就会引入时间与长度变化。
  \end{block}
\end{frame}

\begin{frame}{例子：一次多轮对话的真实链路并不止于模型前向（2）}
  \Core{请求生命周期要覆盖从输入组织到状态回收的完整过程，而不是只盯住模型执行。}
\begin{block}{进入后的主路径}
    入队、调度、prefill、decode、输出组装会在不同对象之间多次转移控制权。
  \end{block}
\end{frame}

\begin{frame}{例子：一次多轮对话的真实链路并不止于模型前向（3）}
  \Core{请求生命周期要覆盖从输入组织到状态回收的完整过程，而不是只盯住模型执行。}
\begin{block}{完成后的尾部动作}
    请求结束不代表系统工作结束，状态释放、统计更新和资源回收同样属于生命周期的一部分。
  \end{block}
\end{frame}

\begin{frame}{例子：chat template 为什么会影响生命周期分析（1）}
  \Core{如果输入在真正进入 engine 前已经发生了长度膨胀，那么后面的阶段观察也会被一并放大。}
\begin{block}{用户看到的是一条自然语言消息}
    但系统里常常先被改写为带角色标记、特殊 token 和 generation prompt 的模板串。
  \end{block}
\end{frame}

\begin{frame}{例子：chat template 为什么会影响生命周期分析（2）}
  \Core{如果输入在真正进入 engine 前已经发生了长度膨胀，那么后面的阶段观察也会被一并放大。}
\begin{block}{直接后果}
    prompt\_len 变化会影响 prefill 代价、KV 初始规模和后续排队行为。
  \end{block}
\end{frame}

\begin{frame}{例子：chat template 为什么会影响生命周期分析（3）}
  \Core{如果输入在真正进入 engine 前已经发生了长度膨胀，那么后面的阶段观察也会被一并放大。}
\begin{block}{教学意义}
    生命周期不是从“用户文本”开始算，而是从“真正进入系统的 token 序列”开始算。
  \end{block}
\end{frame}

\begin{frame}{代码例子：多轮对话在入口处如何被组织（1）}
  \Core{生命周期课程里的代码例子要回答：输入是怎样变成请求对象的。}
\begin{block}{入口片段}
    {\ttfamily\small
    \detokenize{prompts = [tokenizer.apply_chat_template(...)]}\par
    \detokenize{outputs = llm.generate(prompts, sampling_params)}\par
    }
  \end{block}
\end{frame}

\begin{frame}{代码例子：多轮对话在入口处如何被组织（2）}
  \Core{生命周期课程里的代码例子要回答：输入是怎样变成请求对象的。}
\begin{block}{课堂应追问}
    `generate()` 接到的是原始文本、token 序列，还是已经带运行时元数据的请求对象？
  \end{block}
\end{frame}

\begin{frame}{代码例子：多轮对话在入口处如何被组织（3）}
  \Core{生命周期课程里的代码例子要回答：输入是怎样变成请求对象的。}
\begin{block}{系统含义}
    入口处的数据形态，决定了后面 state object 该如何承载长度、状态和 block table。
  \end{block}
\end{frame}

\begin{frame}{代码例子：Sequence 为什么是生命周期里的关键对象（1）}
  \Core{如果不能说清 \texttt{Sequence} 承载了哪些字段，就很难讲明请求是如何跨阶段存活的。}
\begin{block}{课堂应关注的字段}
    {\ttfamily\small
    \detokenize{status}\par
    \detokenize{num_cached_tokens}\par
    \detokenize{block_table}\par
    \detokenize{last_block_num_tokens}\par
    }
  \end{block}
\end{frame}

\begin{frame}{代码例子：Sequence 为什么是生命周期里的关键对象（2）}
  \Core{如果不能说清 \texttt{Sequence} 承载了哪些字段，就很难讲明请求是如何跨阶段存活的。}
\begin{block}{对象视角}
    这些字段把“请求内容、调度状态、KV 元数据”压在了同一个生命周期对象上。
  \end{block}
\end{frame}

\begin{frame}{代码例子：Sequence 为什么是生命周期里的关键对象（3）}
  \Core{如果不能说清 \texttt{Sequence} 承载了哪些字段，就很难讲明请求是如何跨阶段存活的。}
\begin{block}{结论}
    一次请求之所以是系统问题，关键就在于它不是一次函数调用，而是一个持续演化的状态对象。
  \end{block}
\end{frame}

\begin{frame}{例子：prefill 和 decode 为什么应该拆开看（1）}
  \Core{生命周期讲义里最重要的区分，就是 prefill 与 decode 的资源特征不同。}
\begin{block}{prefill}
    一次性处理更长输入，计算更密集，也更容易直接抬高 TTFT。
  \end{block}
\end{frame}

\begin{frame}{例子：prefill 和 decode 为什么应该拆开看（2）}
  \Core{生命周期讲义里最重要的区分，就是 prefill 与 decode 的资源特征不同。}
\begin{block}{decode}
    每轮 token 较少，但会跨更多轮持续存在，更容易与调度和状态驻留绑定在一起。
  \end{block}
\end{frame}

\begin{frame}{例子：prefill 和 decode 为什么应该拆开看（3）}
  \Core{生命周期讲义里最重要的区分，就是 prefill 与 decode 的资源特征不同。}
\begin{block}{误区}
    如果把二者都写成“模型生成阶段”，就会把关键的性能来源混在一起。
  \end{block}
\end{frame}

\begin{frame}{课堂练习：用时序图复述一次请求（1）}
  \Core{课程练习的目的不是画复杂图，而是检验你是否知道控制权何时转移。}
\begin{block}{练习一}
    在时序图中标出：请求进入、prefill 开始、decode 开始、输出完成、状态释放五个节点。
  \end{block}
\end{frame}

\begin{frame}{课堂练习：用时序图复述一次请求（2）}
  \Core{课程练习的目的不是画复杂图，而是检验你是否知道控制权何时转移。}
\begin{block}{练习二}
    说明其中哪一个节点最适合插第一条日志，为什么。
  \end{block}
\end{frame}

\begin{frame}{课堂练习：用时序图复述一次请求（3）}
  \Core{课程练习的目的不是画复杂图，而是检验你是否知道控制权何时转移。}
\begin{block}{练习三}
    用一句话回答：为什么生命周期图是后续调度和 KV 讨论的公共前提？
  \end{block}
\end{frame}


\begin{frame}{本讲总结}
  \Core{理解单个请求如何穿过系统，是分析后续调度、KV 与执行优化问题的基础。}
  \begin{itemize}
    \item 请求链路是系统问题的最小观察单元。
    \item 生命周期分析先于机制设计。
  \end{itemize}
\end{frame}

\LectureThreeSources

\begin{frame}{对应 Tutorial}
  \Core{本讲对应 Tutorial 2。先独立作答，再使用对应 Tutorial 的教师讨论要点核对。}
  \begin{itemize}
    \item 文件：\texttt{\detokenize{build/tutorials/Tutorial_02_请求生命周期与PrefillDecode.pdf}}
    \item 动手运行、日志记录与提交要求统一查看 \texttt{\detokenize{build/experiments/}} 和对应 experiment 源目录下的 \texttt{\detokenize{assignment_spring-2027/}}。
  \end{itemize}
\end{frame}

\end{document}
